\section{The original and zero-base fibre diagrams at pure normal exponents}

Let $S=\mathbb C[[u_1,\ldots,u_d,z_1,\ldots,z_t]]$ and
$T=(u_1,\ldots,u_d)$. Every diagram below is the actual coefficient
diagram of the stated whole ideal. The comparison retains the same
original ideal throughout; equality of the desired diagrams is not an
input.

\begin{theorem}[Literal normal projection preserves pure-normal diagram witnesses]
\label{hir:actual-normal-fibre-diagram-normal-projection}
\lean{Hironaka.HironakaActualNormalFibreDiagramCompare.normalProjection_coefficient,Hironaka.HironakaActualNormalFibreDiagramCompare.pure_normal_diagram_add_tangent_of_normal_extension}
\leanok
Let $\nu:\mathbb C[[z]]\to S$ be literal inclusion of the normal
variables, and suppose a whole ideal $J$ is the extension of its
contraction along $\nu$. For any monomial order $m$ and normal exponent
$\gamma$,
\[
 (0,\gamma)\in\Gamma_m(J+T)
 \quad\Longleftrightarrow\quad
 (0,\gamma)\in\Gamma_m(J).
\]
\end{theorem}
\begin{proof}
Define the actual algebra endomorphism $\tau$ by first killing every
tangent variable and then including the remaining normal series in
$S$. Its coefficient at a pure normal exponent is the original
coefficient, and its other coefficients are zero. The exact retraction
$\tau\nu=\nu$ and the assumed whole normal extension show that
$\tau_*J=J$. Each tangent variable maps to zero, so $\tau_*T=0$.

Suppose a nonzero $F\in J+T$ has leading exponent $(0,\gamma)$.
The coefficient at this exponent is unchanged by $\tau$ and remains
nonzero. Every smaller coefficient of $\tau F$ is either a smaller
coefficient of $F$, hence zero, or has been killed. Thus $\tau F\in J$
has exactly the same genuine leading exponent. This supplies the forward
diagram implication. The reverse implication uses the inclusion
$J\subset J+T$ on the same nonzero witness.
\end{proof}

\begin{theorem}[Whole initial ideals commute with the entire tangent fibre cut]
\label{hir:actual-normal-fibre-diagram-initial-cut}
\lean{Hironaka.HironakaActualNormalFibreDiagramCompare.actualInitialIdeal_add_all_tangent}
\leanok
\uses{hir:actual-initial-coordinate-quotient-ordered-cuts}
Let $K\subset S$. Suppose its actual homogeneous initial ideal is
preserved by every tangent formal partial derivative. Then
\[
 \operatorname{in}(K+T)=\operatorname{in}(K)+T.
\]
\end{theorem}
\begin{proof}
List the first $d$ coordinate indices in their literal order. They are
distinct, since the coordinate-block inclusion is injective. Successive
coordinate quotienting by this list is exactly addition of the span of
all its coordinate variables: induction on the list proves this using
associativity of ideal sums and the span-of-union identity. The proved
successive-coordinate initial-ideal compatibility applies to this same
list, with the given formal tangent differential stability. Substituting
the two explicit sum identities proves the displayed entire ideal
equality.
\end{proof}

\begin{theorem}[Pure-normal original and fibre diagrams agree]
\label{hir:actual-normal-fibre-diagram-original-comparison}
\lean{Hironaka.HironakaActualNormalFibreDiagramCompare.pure_normal_diagram_add_tangent_iff,Hironaka.HironakaActualNormalFibreDiagramCompare.map_originalBaseMaximalIdeal_actualTaylor,Hironaka.HironakaActualNormalFibreDiagramCompare.pure_normal_fibre_diagram_iff}
\leanok
\uses{hir:actual-normal-fibre-diagram-normal-projection,hir:actual-normal-fibre-diagram-initial-cut,hir:actual-whole-initial-normal-extension-of-polynomial-tangent-stability,hir:actual-homogeneous-initial-ideal-diagram-equality,hir:actual-joint-base-fibre-kernel}
Let $I\subset\mathcal O_{\mathbb C^{d+t},0}$ be an original whole
holomorphic germ ideal. Let $K$ be its actual Taylor extension and let
$M$ be the image of the entire tangent-base maximal ideal. Assume the
actual homogeneous initial ideal of $K$ is preserved by all tangent
formal partial derivatives. Then for every normal exponent $\gamma$,
\[
 (0,\gamma)\in\Gamma_{\mathrm{degLex}}
   \bigl(\operatorname{Taylor}_*(I+M)\bigr)
 \quad\Longleftrightarrow\quad
 (0,\gamma)\in\Gamma_{\mathrm{degLex}}(K).
\]
\end{theorem}
\begin{proof}
Formal differential stability implies polynomial differential stability
on the same actual initial ideal, by compatibility of formal and
polynomial partial differentiation. The existing normal-extension
producer therefore makes $\operatorname{in}(K)$ the entire extension
of its actual normal contraction. The first theorem compares its pure
normal diagram with that of $\operatorname{in}(K)+T$. The second
theorem identifies this latter whole ideal with
$\operatorname{in}(K+T)$. The proved actual degree-first diagram equality
between an ideal and its homogeneous initial ideal transfers both sides
back to $K+T$ and $K$.

Finally, the whole tangent-base maximal ideal is generated by its
actual tangent coordinate germs. Taylor sends each such coordinate germ
to precisely the corresponding formal tangent variable. Extension
preserves generating spans and ideal sums, so
$\operatorname{Taylor}_*(I+M)=K+T$ exactly. Applying the established
formal comparison proves the original-germ statement.
\end{proof}

This is equality of membership for every specified pure normal exponent,
not equality of the complete original and fibre diagrams. It supplies the
normal-standard label comparison needed to transport the literal normal
monomial bases; no normal-flatness or Bennett conclusion is asserted here.
